\chapter{La division marchés de la banque}\label{ch:fm:the-bank-markets-division}

Le rapport annuel 2025 d'une grande banque américaine alloue 149.5~milliards de dollars de ses \omterm{def:fm:the-proprietary-market-making-firm:ownfunds}{fonds propres} à sa banque de financement et d'investissement et indique que
l'activité a dégagé 18~\% sur ce montant. Le même rapport explique comment ce chiffre a été retenu~: l'allocation «~intègre les actifs pondérés par les risques de Bâle~III
selon l'approche standardisée et la surcharge des banques d'importance systémique mondiale \dots\ ainsi qu'une simulation de l'épuisement du capital dans un environnement de stress
sévère~». Chaque desk de cette activité répond de sa part des \omterm{def:fm:the-proprietary-market-making-firm:ownfunds}{fonds propres}, et cette part dépend de la règle qui le contraint le plus. Un
desk de repo qui paraît presque gratuit selon une règle est le plus coûteux de l'établissement selon une autre.

\section{Ce que vend une division marchés~: flux, financement et structuration}

La division marchés d'une banque (livre~1, chapitre~2) tire son chiffre d'affaires de trois façons. Les desks de \emph{flux} font du market making pour leurs clients en taux, en change,
en crédit, en actions et en matières premières, et perçoivent la fourchette et la franchise (livre~1, chapitre~2~; livre~9, chapitre~24). Les desks de \emph{financement} prêtent des liquidités
et des titres à leurs clients, en repo, en prêt de titres et en services de prime brokerage, et perçoivent un spread sur le bilan qu'ils prêtent. Les desks de \emph{structuration}
vendent des dérivés sur mesure et des produits structurés et perçoivent une marge qui rémunère le risque qu'ils portent (livre~5, chapitre~19).

Les trois activités n'utilisent pas les ressources de la banque dans des proportions très différentes. Un desk de flux fait tourner son inventaire et consomme du capital de risque pour le
risque de marché qu'il détient~; un desk de financement détient de gros actifs peu risqués au bilan pour un spread mince~; un desk de structuration porte un risque complexe et de longue
échéance, très consommateur de capital par dollar de notionnel. Ce que vend la division se mesure donc non en chiffre d'affaires, mais en chiffre d'affaires par
unité de la contrainte que consomme chaque activité.

\begin{omfigure}
\begin{tikzpicture}[font=\footnotesize,>=stealth,box/.style={draw,rounded corners,minimum height=0.62cm,minimum width=2.5cm,align=center,
  fill=omDef!10},res/.style={draw,rounded corners,minimum height=0.62cm,minimum width=2.6cm,align=center,fill=omProp!12}]
\node[box] (flow) at (0,2.6) {desks de flux\\{\scriptsize fourchette, franchise}};
\node[box] (fin) at (0,1.3) {desks de financement\\{\scriptsize repo, prime}};
\node[box] (str) at (0,0) {structuration\\{\scriptsize risque sur mesure}};
\node[res] (rwa) at (5.4,2.6) {actifs pondérés par les risques};
\node[res] (le) at (5.4,1.3) {exposition au levier};
\node[res] (st) at (5.4,0) {perte de stress};
\draw[->,thick] (flow) -- (rwa); \draw[->,thin,gray] (flow) -- (le);
\draw[->,thick] (fin) -- (le); \draw[->,thin,gray] (fin) -- (rwa);
\draw[->,thick] (str) -- (st); \draw[->,thick] (str) -- (rwa);
\node[draw,rounded corners,fill=omThm!10,minimum width=2.4cm,minimum height=0.62cm,align=center] (eq) at (9.6,1.3) {fonds propres du groupe\\{\scriptsize le besoin le plus élevé}};
\draw[->] (rwa) -- (eq); \draw[->] (le) -- (eq); \draw[->] (st) -- (eq);
\end{tikzpicture}

\omcaption{Ce que consomme chaque activité d'une division marchés. Les desks de flux et de structuration utilisent des actifs pondérés par les risques et de la capacité de stress~; les desks de financement
utilisent du bilan, mesuré par l'exposition au levier. Le groupe détient des \omterm{def:fm:the-proprietary-market-making-firm:ownfunds}{fonds propres} face au plus grand de ses besoins~; les flèches épaisses marquent la
consommation principale de chaque activité.}\label{fig:fm:the-bank-markets-division:lines}
\end{omfigure}

\section{La pile de capital~: fonds propres de base, levier et coussins de stress}

\begin{definition}[Fonds propres de base de catégorie 1]\label{def:fm:the-bank-markets-division:cet1}
Les \emph{fonds propres de base de catégorie 1}\index{fonds propres de base de catégorie 1} (CET1) d'une banque sont les capitaux propres des actionnaires ordinaires, les bénéfices non
distribués et réserves, diminués des déductions (écarts d'acquisition et autres immobilisations incorporelles, certains actifs d'impôt différé)~; ce sont les \omterm{def:fm:the-proprietary-market-making-firm:ownfunds}{fonds propres} qui absorbent les pertes en premier
tant que la banque est en activité. Son ratio CET1 est le rapport des \omterm{def:fm:the-proprietary-market-making-firm:ownfunds}{fonds propres} CET1 aux actifs pondérés par les risques.
\end{definition}

\begin{definition}[Ratio de levier, mesure de l'exposition au levier]\label{def:fm:the-bank-markets-division:lr}
Le \emph{ratio de levier}\index{ratio de levier} (leverage ratio) d'une banque est le rapport de son capital de catégorie 1 à sa
\emph{mesure de l'exposition au levier}\index{mesure de l'exposition au levier} (leverage exposure measure)~: ses actifs au bilan, diminués des montants déduits du capital de catégorie 1, augmentés de certaines expositions hors bilan (exposition
potentielle des dérivés, financement de titres, engagements), sans pondération par le risque.
\end{definition}

\begin{definition}[Coussin de fonds propres de stress, surcharge G-SIB]\label{def:fm:the-bank-markets-division:buffers}
Un \emph{coussin de fonds propres de stress}\index{coussin de fonds propres de stress} (stress capital buffer) est un coussin de capital fixé à partir de la perte de capital qu'une banque subirait dans le scénario de stress sévère
du superviseur, et ajouté à son ratio minimal. Une \emph{surcharge G-SIB}\index{surcharge G-SIB} est un coussin de capital supplémentaire exigé
d'une banque d'importance systémique mondiale, fixée à partir d'un score mesurant sa taille, son interconnexion, sa complexité, son activité transfrontalière et
sa substituabilité.
\end{definition}

\begin{dated}{2026-09}{Pile de capital d'une banque à la fin de 2025}\label{dat:fm:the-bank-markets-division:jpm}
JPMorgan Chase \& Co.\ (formulaire 10-K pour 2025), approche standardisée~: \omterm{def:fm:the-proprietary-market-making-firm:ownfunds}{fonds propres} CET1 de 288.5~milliards de dollars pour des RWA de 1\,981.7~milliards de dollars, un ratio CET1 de
14.6~\% face à une exigence de 11.5~\% coussins compris (un minimum de 4.5~\%, un coussin de capital de stress de 2.5~\%, le plancher, en vigueur
jusqu'au 30 septembre 2027, et une surcharge GSIB de 4.5~\% selon la méthode~2~; 2.5~\% selon la méthode~1). \omterm{def:fm:the-bank-markets-division:lr}{Ratio de levier} supplémentaire de 5.8~\% sur une exposition totale au levier
de 5\,302.0~milliards de dollars, face à un minimum de 3.0~\% plus un coussin de 2.0~\%. En novembre 2025, la Réserve fédérale, l'OCC et la FDIC ont remplacé les
coussins SLR renforcés et statiques des plus grandes banques par un coussin égal à 50~\% de la surcharge GSIB selon la méthode~1 de la société holding (plafonné à 1~\% pour ses
filiales bancaires). \omterm{def:fm:the-bank-markets-division:alloc}{Fonds propres alloués} à la banque de financement et d'investissement~: 149.5~milliards de dollars à la fin de 2025, 166.5~milliards de dollars à partir du 1\up{er} janvier 2026~;
le résultat net de l'activité en 2025 s'élève à 27\,761~millions de dollars et son ROE à 18~\%.
\end{dated}

La pile fixe deux exigences qui ne mordent pas au même endroit (\cref{fig:fm:the-bank-markets-division:stack}). Pour la banque de l'encadré daté,
les règles fondées sur les risques exigent un capital de catégorie 1 de 13.0~\% des RWA, soit 257.6~milliards de dollars~; la règle de levier exige 5.0~\% de son exposition au levier, soit 265.1
milliards de dollars. Au niveau de l'ensemble de la banque, c'est la règle de levier qui contraint, de peu~; son exposition au levier vaut 2.7 fois ses RWA. Une division marchés
dont l'activité de financement porte son exposition au levier à quatre fois ses RWA se heurte à une règle de levier qui contraint nettement.

\begin{omfigure}
\begin{tikzpicture}
\begin{axis}[width=10.5cm,height=5cm,xbar stacked,bar width=14pt,xmin=0,xmax=16,ytick={0,1},
  yticklabels={{CET1 / actifs pondérés par les risques},{Tier 1 / exposition au levier}},y dir=reverse,ymin=-0.6,ymax=1.6,
  xlabel={\small en pourcentage},tick label style={font=\footnotesize},table/col sep=comma,grid=major,grid style={gray!20},
  legend cell align=left,legend style={font=\scriptsize,at={(0.5,-0.32)},anchor=north,/tikz/every even column/.append style={column sep=8pt}},legend columns=4]
\addplot[fill=omDef!55,draw=omDef] table[y=k,x=minimum]{figdata/desk/05-the-bank-markets-division/stack.csv};
\addplot[fill=omProp!55,draw=omProp] table[y=k,x=buffer]{figdata/desk/05-the-bank-markets-division/stack.csv};
\addplot[fill=omThm!50,draw=omThm] table[y=k,x=gsib]{figdata/desk/05-the-bank-markets-division/stack.csv};
\addlegendimage{only marks,mark=diamond*,mark size=3.5pt,black}
\legend{minimum,{coussin de stress ou de levier},surcharge GSIB,{ratio effectif}}
\end{axis}
\begin{axis}[width=10.5cm,height=5cm,xmin=0,xmax=16,y dir=reverse,ymin=-0.6,ymax=1.6,hide axis,table/col sep=comma]
\addplot[only marks,mark=diamond*,mark size=3.5pt,black] table[y=k,x=actual]{figdata/desk/05-the-bank-markets-division/stack.csv};
\end{axis}
\end{tikzpicture}

\omcaption{Les deux exigences d'une banque à la fin de 2025~: CET1 sur actifs pondérés par les risques (un minimum de 4.5~\%, un coussin de capital de stress de 2.5~\% et une
surcharge GSIB de 4.5~\%, soit 11.5~\%, face à 14.6~\% détenus) et Tier 1 sur exposition au levier (3.0~\% plus un coussin de 2.0~\%, face à 5.8~\%). Données~: JPMorgan
Chase \& Co., formulaire 10-K pour 2025.}\label{fig:fm:the-bank-markets-division:stack}
\end{omfigure}

\section{Rentabilité du capital par desk}

\begin{definition}[Fonds propres alloués, contrainte de capital contraignante]\label{def:fm:the-bank-markets-division:alloc}
Les \emph{fonds propres alloués}\index{fonds propres alloués} (allocated equity) d'un desk sont la part des \omterm{def:fm:the-proprietary-market-making-firm:ownfunds}{fonds propres} du groupe qui lui est imputée, calculée à partir de ce qu'il consomme de
chaque contrainte à laquelle le groupe est soumis. La \emph{contrainte de capital contraignante}\index{contrainte de capital contraignante} (binding capital constraint) d'un desk, ou du groupe, est la
contrainte qui exige le plus de \omterm{def:fm:the-proprietary-market-making-firm:ownfunds}{fonds propres} pour lui~: actifs pondérés par les risques, exposition au levier ou perte de stress.
\end{definition}

La rentabilité sur \omterm{def:fm:the-bank-markets-division:alloc}{fonds propres alloués} d'un desk est son résultat après impôt rapporté à ses \omterm{def:fm:the-bank-markets-division:alloc}{fonds propres alloués}, version au niveau du desk de la \omterm{def:fm:the-economics-of-a-trading-firm:roe}{rentabilité des fonds propres}
(ROE) du chapitre~1, et elle se compare au taux de rendement minimal exigé de la banque (livre~6, chapitre~19). La division du chapitre (illustrative) compte six desks.
Avec des \omterm{def:fm:the-proprietary-market-making-firm:ownfunds}{fonds propres} égaux à 13~\% des actifs pondérés par les risques, 5~\% de l'exposition au levier, ou la perte de stress du desk, un taux d'impôt de 25~\% et les chiffres des desks figurant dans le
module (\texttt{fm\_bank.DESKS}), leurs rentabilités sont~:

\begin{center}
\small
\begin{tabular}{@{}lrrrl@{}}
\toprule
desk & clé RWA & clé levier & contrainte liante (la plus forte) & contrainte active\\
\midrule
flux de taux & 15.4~\% & 9.6~\% & 9.6~\% & levier\\
repo & 53.8~\% & 5.2~\% & 5.2~\% & levier\\
services de prime brokerage & 23.1~\% & 13.5~\% & 13.5~\% & levier\\
dérivés actions & 18.1~\% & 27.5~\% & 18.1~\% & RWA\\
FX & 28.8~\% & 28.1~\% & 28.1~\% & levier\\
crédit & 10.4~\% & 22.5~\% & 10.4~\% & RWA\\
\bottomrule
\end{tabular}
\end{center}

\begin{omfigure}
\begin{tikzpicture}
\begin{axis}[width=11.5cm,height=5.4cm,ybar,bar width=8pt,ymin=0,ymax=60,ylabel={\small rentabilité sur fonds propres alloués (\%)},
  xtick={0,1,2,3,4,5},xticklabels={{flux de taux},repo,{services de prime brokerage},{dérivés actions},FX,crédit},
  xticklabel style={font=\scriptsize,text width=1.5cm,align=center},tick label style={font=\footnotesize},table/col sep=comma,
  enlarge x limits=0.1,grid=major,grid style={gray!20},legend cell align=left,
  legend style={font=\scriptsize,at={(0.5,-0.32)},anchor=north,/tikz/every even column/.append style={column sep=8pt}},legend columns=3]
\addplot[fill=omDef!55,draw=omDef] table[x=k,y=rwa]{figdata/desk/05-the-bank-markets-division/roae.csv};
\addplot[fill=omThm!50,draw=omThm] table[x=k,y=leverage]{figdata/desk/05-the-bank-markets-division/roae.csv};
\addplot[fill=gray!40,draw=gray] table[x=k,y=binding]{figdata/desk/05-the-bank-markets-division/roae.csv};
\draw[omProp,dashed,thick] (axis cs:-0.6,12) -- (axis cs:5.6,12);
\legend{fonds propres par RWA,fonds propres par exposition au levier,fonds propres selon la contrainte la plus forte}
\end{axis}
\end{tikzpicture}

\omcaption{Les rentabilités après impôt sur \omterm{def:fm:the-bank-markets-division:alloc}{fonds propres alloués} des six desks, selon les actifs pondérés par les risques, selon l'exposition au levier, et selon celle qui exige le plus~;
la ligne en pointillés est le taux de rendement minimal exigé de 12~\%. Desks illustratifs~; niveaux d'exigence à l'échelle de l'encadré daté. Données~:
\texttt{fm\_bank.roae\_table}.}\label{fig:fm:the-bank-markets-division:roae}
\end{omfigure}

Le desk repo est la meilleure activité de la division en actifs pondérés par les risques et la pire en levier. Ses actifs sont des prêts garantis à court terme, assortis de faibles
pondérations de risque, et ils remplissent le bilan~: 400~milliards de dollars d'exposition au levier pour 15~milliards de dollars de RWA. Les dérivés actions et le crédit
sont à l'opposé~: lourds en risque et légers en bilan, ils sont contraints par les RWA.

\begin{proposition}[Une allocation selon le maximum propre à chaque desk surestime les fonds propres]\label{prop:fm:the-bank-markets-division:max}
Soit le desk $i$ qui a besoin de $a_i$ de \omterm{def:fm:the-proprietary-market-making-firm:ownfunds}{fonds propres} sous une contrainte et de $b_i$ sous une autre, et soit $E=\max(\sum_ia_i,\sum_ib_i)$ les \omterm{def:fm:the-proprietary-market-making-firm:ownfunds}{fonds propres} détenus par le groupe. Alors
$\sum_i\max(a_i,b_i)\ge E$, avec égalité seulement si la même contrainte est contraignante pour chaque desk.
\end{proposition}

\begin{proof}
$\max(a_i,b_i)\ge a_i$ et $\ge b_i$ pour tout $i$~; on additionne chaque inégalité. L'égalité dans les deux sommes exige $\max(a_i,b_i)=a_i$ pour tout $i$ ou
$=b_i$ pour tout $i$.
\end{proof}

Les desks du chapitre ont besoin de 35.1~milliards de dollars de \omterm{def:fm:the-proprietary-market-making-firm:ownfunds}{fonds propres} selon les RWA et de 55.5~milliards selon le levier~; le groupe en détient 55.5~milliards de dollars, mais les maxima propres aux desks
totalisent 62.1~milliards de dollars. Si chacun est imputé selon sa propre contrainte contraignante, la division afficherait une rentabilité de 11.4~\% sur des \omterm{def:fm:the-proprietary-market-making-firm:ownfunds}{fonds propres} qu'elle n'a
pas. Une allocation cohérente ramène une clé à la contrainte contraignante du groupe, ou mélange les clés avec des pondérations dont la somme égale les \omterm{def:fm:the-proprietary-market-making-firm:ownfunds}{fonds propres} du groupe,
comme le fait la banque de l'encadré daté lorsqu'elle combine les RWA, sa \omterm{def:fm:the-bank-markets-division:buffers}{surcharge G-SIB} et une simulation de stress.

\section{Le bilan comme ressource rare}

\begin{definition}[Coût du bilan]\label{def:fm:the-bank-markets-division:charge}
Un \emph{coût du bilan}\index{coût du bilan} (balance-sheet charge) est un prix interne par unité de bilan, généralement d'exposition au levier, facturé
à un desk pour ce qu'il consomme~: le coût des \omterm{def:fm:the-proprietary-market-making-firm:ownfunds}{fonds propres} que l'exigence de levier impose au groupe de détenir en contrepartie.
\end{definition}

\begin{method}[Fixer un coût du bilan]\label{met:fm:the-bank-markets-division:charge}
\begin{enumerate}
\item Prendre les \omterm{def:fm:the-proprietary-market-making-firm:ownfunds}{fonds propres} que l'exigence de levier impose par unité d'exposition au levier (5~\% ici, ou la cible de la banque, coussins compris).
\item Multiplier par le taux de rendement minimal exigé~: à 12~\%, le coût est $0.12\times0.05=0.6~\%$ par an d'exposition au levier.
\item Déduire le coût du résultat avant impôt de chaque desk, à côté du taux de cession interne (funds-transfer pricing, livre~6, chapitre~24) pour son financement.
\item Un desk dont le résultat après \omterm{def:fm:the-bank-markets-division:charge}{coût du bilan} est négatif ne rémunère pas son coût du capital sur le bilan qu'il utilise.
\end{enumerate}
\end{method}

À 0.6~\% par an, l'exposition au levier de 400~milliards de dollars du desk repo coûte 2.4~milliards de dollars, et son résultat après impôt de 1.05~milliard de dollars devient une perte de
0.75~milliard de dollars. Le seuil de rentabilité est atteint à un coût de 0.35~\% (35 points de base) par an. Le desk de services de prime brokerage gagne 0.45~milliard de dollars après ce coût, et les
dérivés actions 1.11~milliard de dollars. Le \omterm{def:fm:the-bank-markets-division:charge}{coût du bilan} ne dit pas que le repo doit fermer~; il dit que chaque dollar de bilan qu'il utilise doit être
facturé à 60 points de base par an, et que l'activité n'est viable qu'à un spread supérieur à ce niveau.

\omcode{code/firm/bankdesk/firm_bankdesk.py}{45}{79}{Fonds propres par chaque contrainte et par la contrainte contraignante, rentabilités correspondantes, coût du bilan,
et le mix des desks sous les contraintes du groupe.}\label{lst:fm:the-bank-markets-division:bankdesk}

Les \omterm{def:fm:the-proprietary-market-making-firm:ownfunds}{fonds propres} du groupe étant fixés aux 55.5~milliards de dollars que son exigence de levier impose aujourd'hui, le mix de desks qui maximise le résultat après impôt de la division
(un programme linéaire, \texttt{firm.bankdesk.optimal\_mix}, chaque desk pouvant varier entre zéro et 1.5 fois sa taille) réduit le repo à environ un
dixième de sa taille et porte chaque autre desk à la limite (\cref{fig:fm:the-bank-markets-division:mix})~: le résultat passe de 7.05~milliards de dollars à 9.12
milliards de dollars. Le programme suppose que chaque desk peut croître sans réduire ses marges et que les desks sont indépendants. La seconde hypothèse est celle
que brise une franchise.

\begin{omfigure}
\begin{tikzpicture}
\begin{axis}[width=11.5cm,height=5cm,ybar,bar width=10pt,ymin=0,ymax=1.7,ylabel={\small échelle du desk (aujourd'hui $=1$)},
  xtick={0,1,2,3,4,5},xticklabels={{flux de taux},repo,{services de prime brokerage},{dérivés actions},FX,crédit},
  xticklabel style={font=\scriptsize,text width=1.5cm,align=center},tick label style={font=\footnotesize},table/col sep=comma,
  enlarge x limits=0.1,grid=major,grid style={gray!20},legend cell align=left,
  legend style={font=\scriptsize,at={(0.5,-0.32)},anchor=north,/tikz/every even column/.append style={column sep=8pt}},legend columns=2]
\addplot[fill=gray!40,draw=gray] table[x=k,y=today]{figdata/desk/05-the-bank-markets-division/mix.csv};
\addplot[fill=omDef!55,draw=omDef] table[x=k,y=optimal]{figdata/desk/05-the-bank-markets-division/mix.csv};
\legend{aujourd'hui,mix maximisant le résultat sous les deux contraintes du groupe}
\end{axis}
\end{tikzpicture}

\omcaption{Le mix de desks qui maximise le résultat après impôt de la division avec les \omterm{def:fm:the-proprietary-market-making-firm:ownfunds}{fonds propres} actuels, lorsque chaque desk peut se réduire à zéro ou croître de moitié et
que les marges ne changent pas~: le desk repo, gros consommateur de bilan, cède son exposition au levier aux autres. Données~: \texttt{fm\_bank.mix},
\texttt{firm.bankdesk.optimal\_mix}.}\label{fig:fm:the-bank-markets-division:mix}
\end{omfigure}

\section{La franchise clientèle}

Un client qui finance ses positions auprès d'une banque compense, emprunte et négocie aussi avec elle. La franchise des desks de flux (livre~1, chapitre~2)
repose en partie sur le bilan des desks de financement~: les clients des services de prime brokerage choisissent la banque qui leur prête, et négocient avec elle~; le central
risk book (livre~9, chapitre~25) compense les flux de tous. Un desk qui paraît détruire de la valeur pris isolément peut être ce qui retient les clients dans la
maison.

\begin{example}[Réduction du repo]\label{ex:fm:the-bank-markets-division:cut}
La division réduit de moitié son desk repo, et les services de prime brokerage perdent 20~\% de leur produit, à coûts inchangés, car certains clients financent
ailleurs leurs positions. Le résultat après impôt passe de 7.05~milliards de dollars à 6.08~milliards de dollars~; les \omterm{def:fm:the-proprietary-market-making-firm:ownfunds}{fonds propres} dont ont besoin les desks passent de 62.1~milliards de dollars à 50.1~milliards de dollars (chacun
par sa \omterm{def:fm:the-bank-markets-division:alloc}{contrainte de capital contraignante})~; la rentabilité de ces \omterm{def:fm:the-proprietary-market-making-firm:ownfunds}{fonds propres} passe de 11.4~\% à 12.1~\%. Savoir si la réduction est judicieuse dépend de ce que rapportent les \omterm{def:fm:the-proprietary-market-making-firm:ownfunds}{fonds propres} libérés~:
restitués aux actionnaires, ou donnés aux desks que le programme linéaire voulait faire croître.
\end{example}

\begin{remark}[Pourquoi les banques se spécialisent]\label{rem:fm:the-bank-markets-division:special}
Comme la \omterm{def:fm:the-bank-markets-division:alloc}{contrainte de capital contraignante} de chaque banque diffère, une même activité rapporte des rendements différents selon les banques~: une banque contrainte par le levier tarifie le repo
cher et les dérivés sur actions à bon marché~; une banque contrainte par les actifs pondérés par les risques fait l'inverse. Les règles sont les mêmes pour toutes les grandes banques~; leurs
bilans ne le sont pas, et les activités que gagne chaque banque suivent ses contraintes (le livre~17 en dresse le bilan).
\end{remark}

\section{Tutoriel~: six desks et trois clés}\label{tut:fm:the-bank-markets-division:tut}

\begin{tutorial}
\textbf{Objectif.} Répartir les \omterm{def:fm:the-proprietary-market-making-firm:ownfunds}{fonds propres} d'une division selon chaque contrainte, calculer le rendement de chaque desk, tarifer le bilan et trouver le meilleur mix.
\textbf{État final~:} le tableau des rendements et les \cref{fig:fm:the-bank-markets-division:roae,fig:fm:the-bank-markets-division:mix}.
\begin{tutsteps}
\item \textbf{La banque.} \texttt{fm\_bank.jpm\_requirements()} lit les chiffres publiés et calcule le Tier 1 qu'exige chaque règle~:
257.6~milliards de dollars au titre des RWA, 265.1~milliards de dollars au titre de la \omterm{def:fm:the-bank-markets-division:lr}{mesure de l'exposition au levier}.
\item \textbf{Clés.} \texttt{firm.bankdesk.allocate(DESKS, KEYS)} donne les \omterm{def:fm:the-proprietary-market-making-firm:ownfunds}{fonds propres} de chaque desk selon les RWA, le levier et le stress, et retient le plus élevé~;
\texttt{roae} divise le résultat après impôt par chacun de ces montants (voir le \cref{lst:fm:the-bank-markets-division:bankdesk}).
\item \textbf{Le \omterm{def:fm:the-bank-markets-division:charge}{coût du bilan}.} \texttt{fm\_bank.charged\_profits()} déduit 0.6~\% de l'exposition au levier~; \texttt{breakeven\_charge} calcule pour
chaque desk son seuil de rentabilité.
\item \textbf{Le mix.} \texttt{fm\_bank.mix()} résout le programme linéaire~; \texttt{franchise()} lance l'\cref{ex:fm:the-bank-markets-division:cut}.
\end{tutsteps}
\end{tutorial}

\textbf{Pour aller plus loin.} Relever l'exigence de levier du montant du nouveau coussin eSLR et refaire la répartition~; ajouter une contrainte selon laquelle les services de prime brokerage
ne peuvent pas dépasser la taille du repo, un lien de franchise grossier, puis relancer le programme.

\section{Construction~: le modèle de capital de la division}\label{bld:fm:the-bank-markets-division:bankdesk}

\begin{build}
\bfield{Rôle} La partie capital de la division de marchés d'une banque~: ce que coûte chaque desk en \omterm{def:fm:the-proprietary-market-making-firm:ownfunds}{fonds propres} selon chaque règle, ce qu'il rapporte sur ces \omterm{def:fm:the-proprietary-market-making-firm:ownfunds}{fonds propres}, et ce que vaut le
bilan.
\bfield{Interface} \texttt{firm.bankdesk}~: \texttt{Desk(name, revenue, cost, rwa, le, stress)}~; \texttt{Keys(k\_rwa, k\_le, tax)}~;
\texttt{allocate}, \texttt{roae}, \texttt{binding}~; \texttt{balance\_sheet\_charge}, \texttt{breakeven\_charge}~;
\texttt{optimal\_mix(desks, equity, keys, x\_max)}.
\bfield{Règles} Chaque exigence est linéaire dans la taille du desk~; l'allocation contraignante est signalée par un avertissement, car sa somme dépasse les \omterm{def:fm:the-proprietary-market-making-firm:ownfunds}{fonds propres}
du groupe (la \cref{prop:fm:the-bank-markets-division:max})~; les contraintes du programme linéaire sont celles du groupe, jamais celles des desks.
\bfield{Tests d'acceptation} \texttt{code/firm/bankdesk/tests/}~: \omterm{def:fm:the-proprietary-market-making-firm:ownfunds}{fonds propres} par clé et par contrainte contraignante, sur un exemple à deux desks~; le coût de seuil de rentabilité annule le
résultat~; le mix optimal respecte les deux contraintes et favorise le desk le moins gourmand en bilan.
\bfield{Pour aller plus loin} Une contrainte de stress dans le programme~; une allocation mixte dont la somme égale les \omterm{def:fm:the-proprietary-market-making-firm:ownfunds}{fonds propres} du groupe~; le capital FRTB issu du
livre~6, \texttt{firm.frtb}, comme entrée RWA pour les desks de trading.
\end{build}

\begin{omsources}
\item JPMorgan Chase \& Co., Form 10-K for 2025 (résultats par segment, fonds propres par ligne d'activité, gestion du risque de capital).
\item The Goldman Sachs Group, Inc., Form 10-K for 2025 (résultats par segment).
\item One Quant Book 6, chapitres 19, 23 et 24 (taux de rendement minimal exigé, capital du portefeuille de négociation, taux de cession interne).
\end{omsources}

\section{Exercices}

\begin{exercise}[$\star$]\label{exo:fm:the-bank-markets-division:1}
La banque de l'encadré daté disposait de \omterm{def:fm:the-proprietary-market-making-firm:ownfunds}{fonds propres} de Tier 1 de 307.6~milliards de dollars et d'une exposition totale au levier de 5\,302.0~milliards de dollars. Calculer son \omterm{def:fm:the-bank-markets-division:lr}{ratio de
levier} supplémentaire et le Tier 1 que couvre son exigence de 5.0~\%.
\end{exercise}

\begin{exercise}[$\star$]\label{exo:fm:the-bank-markets-division:2}
Diviser le résultat net de 2025 de la banque commerciale et d'investissement, soit 27\,761~millions de dollars, par ses 149.5~milliards de dollars de \omterm{def:fm:the-bank-markets-division:alloc}{fonds propres alloués}, et comparer avec le
ROE publié de 18~\%. Pourquoi peuvent-ils différer~?
\end{exercise}

\begin{exercise}[$\star$]\label{exo:fm:the-bank-markets-division:3}
Le desk repo dégage 2.0~milliards de dollars de produit pour 0.6~milliard de dollars de coûts, utilise 15~milliards de dollars de RWA et 400~milliards de dollars d'exposition au levier. Calculer son
rendement après impôt des \omterm{def:fm:the-proprietary-market-making-firm:ownfunds}{fonds propres}, par clé, à 13~\% et 5~\%, avec un taux d'imposition de 25~\%.
\end{exercise}

\begin{exercise}[$\star\star$]\label{exo:fm:the-bank-markets-division:4}
Avec un taux de rendement minimal exigé de 12~\% et une exigence de levier de 5~\%, quel est le \omterm{def:fm:the-bank-markets-division:charge}{coût du bilan}, et à quel coût le desk repo atteint-il son seuil de rentabilité~?
\end{exercise}

\begin{exercise}[$\star\star$]\label{exo:fm:the-bank-markets-division:5}
Montrer que les allocations maximales propres des desks totalisent 62.1~milliards de dollars alors que le groupe détient 55.5~milliards de dollars, et expliquer l'écart à l'aide de la
\cref{prop:fm:the-bank-markets-division:max}.
\end{exercise}

\begin{exercise}[$\star\star$]\label{exo:fm:the-bank-markets-division:6}
Dans la \cref{fig:fm:the-bank-markets-division:roae}, quels desks franchissent le taux de rendement minimal de 12~\% avec l'allocation contraignante, et lesquels seulement avec la clé RWA~?
\end{exercise}

\begin{exercise}[$\star\star\star$]\label{exo:fm:the-bank-markets-division:7}
\emph{Programmation.} Relancer \texttt{fm\_bank.franchise} en supposant que les services de prime brokerage perdent 40~\% de leur produit au lieu de 20~\%. Que deviennent le résultat et le
rendement des \omterm{def:fm:the-proprietary-market-making-firm:ownfunds}{fonds propres} des desks~?
\end{exercise}

\begin{exercise}[$\star\star\star$]\label{exo:fm:the-bank-markets-division:8}
\emph{Trouver la faille.} «~Le programme linéaire indique que le repo doit diminuer de 89~\%. Faisons-le dès ce trimestre.~»
\end{exercise}

\section{Problème~: le desk qui n'utilise aucun capital}

\begin{problem}[{Problème du week-end --- le desk qui n'utilise aucun capital}]\label{pb:fm:the-bank-markets-division:1}
Le responsable d'une division de marchés entend dire que le desk repo «~n'utilise aucun capital~» et qu'il en tire 54~\%. Il s'agit de vérifier.

\medskip\noindent\textbf{Partie I --- Les règles.}
\begin{enumerate}
\item Définir les \omterm{def:fm:the-bank-markets-division:cet1}{fonds propres de base de catégorie 1} (CET1), le \omterm{def:fm:the-bank-markets-division:lr}{ratio de levier} et la \omterm{def:fm:the-bank-markets-division:lr}{mesure de l'exposition au levier}.
\item Décomposer l'exigence de CET1 de 11.5~\% publiée par la banque.
\item Calculer le Tier 1 que couvrent ses exigences fondées sur les risques et son exigence de levier, et dire laquelle est contraignante.
\item Qu'a modifié la règle de novembre 2025 dans les coussins de levier des plus grandes banques~?
\item Comment la banque publiée répartit-elle ses \omterm{def:fm:the-proprietary-market-making-firm:ownfunds}{fonds propres} entre ses lignes d'activité~?
\end{enumerate}

\medskip\noindent\textbf{Partie II --- Les desks.}
\begin{enumerate}[resume]
\item Définir les \omterm{def:fm:the-bank-markets-division:alloc}{fonds propres alloués} et la \omterm{def:fm:the-bank-markets-division:alloc}{contrainte de capital contraignante}.
\item Donner le rendement du desk repo selon les RWA, selon le levier et selon la \omterm{def:fm:the-bank-markets-division:alloc}{contrainte de capital contraignante}.
\item Quels desks sont contraints par les RWA, et pourquoi~?
\item Énoncer et démontrer la \cref{prop:fm:the-bank-markets-division:max}.
\item Donner les \omterm{def:fm:the-proprietary-market-making-firm:ownfunds}{fonds propres} selon les RWA, selon le levier et selon les maxima au niveau des desks, ainsi que le rendement de la division sur ce dernier.
\end{enumerate}

\medskip\noindent\textbf{Partie III --- Le bilan.}
\begin{enumerate}[resume]
\item Définir un \omterm{def:fm:the-bank-markets-division:charge}{coût du bilan} et le calculer pour un taux de rendement minimal exigé de 12~\%.
\item Donner le résultat net après impôt de chaque desk après déduction de ce coût.
\item Pour quel coût le desk repo atteint-il le seuil de rentabilité~?
\item Que fait le programme linéaire des \omterm{def:fm:the-proprietary-market-making-firm:ownfunds}{fonds propres} actuels, et quel profit atteint-il~?
\item Quelles sont les deux hypothèses du programme les plus fragiles~?
\end{enumerate}

\medskip\noindent\textbf{Partie IV --- La franchise.}
\begin{enumerate}[resume]
\item Que deviennent le profit, les \omterm{def:fm:the-proprietary-market-making-firm:ownfunds}{fonds propres} et le rendement lorsque le repo est réduit de moitié et que les services de prime brokerage perdent 20~\% de leur chiffre d'affaires~?
\item Pourquoi une banque contrainte par le levier et une banque contrainte par les RWA peuvent-elles exercer des activités différentes~?
\item Que demanderiez-vous au desk de services de prime brokerage avant de réduire le repo~?
\item Énoncer le \emph{résultat clé}~: le rendement du desk repo selon les clés RWA et levier, et son \omterm{def:fm:the-bank-markets-division:charge}{coût du bilan} au seuil de rentabilité.
\item En deux phrases, répondre à l'affirmation selon laquelle le repo «~n'utilise aucun capital~».
\end{enumerate}
\end{problem}

\section{Questions d'entretien}

\begin{interviewq}[$\star$ \iqroles{banque}]\label{iq:fm:the-bank-markets-division:1}
Quelle est la différence entre un ratio de \omterm{def:fm:the-proprietary-market-making-firm:ownfunds}{fonds propres} pondéré par les risques et un \omterm{def:fm:the-bank-markets-division:lr}{ratio de levier}, et pourquoi les régulateurs imposent-ils les deux~?
\end{interviewq}

\begin{interviewq}[$\star$ \iqroles{banque, trader}]\label{iq:fm:the-bank-markets-division:2}
Pourquoi un desk repo est-il léger en capital selon les actifs pondérés par les risques et lourd selon le levier~?
\end{interviewq}

\begin{interviewq}[$\star\star$ \iqroles{banque}]\label{iq:fm:the-bank-markets-division:3}
Un desk gagne 1~milliard de dollars après impôt et utilise 200~milliards de dollars d'exposition au levier. Pour une exigence de levier de 5~\% et un taux de rendement minimal exigé de 12~\%, crée-t-il
de la valeur~?
\end{interviewq}

\begin{interviewq}[$\star\star$ \iqroles{banque, risque}]\label{iq:fm:the-bank-markets-division:4}
Pourquoi allouer à chaque desk des \omterm{def:fm:the-proprietary-market-making-firm:ownfunds}{fonds propres} selon sa propre \omterm{def:fm:the-bank-markets-division:alloc}{contrainte de capital contraignante} peut-il surestimer les \omterm{def:fm:the-proprietary-market-making-firm:ownfunds}{fonds propres} dont le groupe a besoin~?
\end{interviewq}

\begin{interviewq}[$\star\star$ \iqroles{banque}]\label{iq:fm:the-bank-markets-division:5}
Deux banques sont soumises aux mêmes règles. Pourquoi l'une pourrait-elle gérer une grande activité de repo et l'autre une grande activité de produits dérivés actions~?
\end{interviewq}

\begin{interviewq}[$\star\star\star$ \iqroles{banque, chercheur}]\label{iq:fm:the-bank-markets-division:6}
On vous demande de fixer la répartition des desks pour l'année prochaine avec un montant de \omterm{def:fm:the-proprietary-market-making-firm:ownfunds}{fonds propres} donné. Qu'ajouteriez-vous à un programme linéaire
avant de faire confiance à sa réponse~?
\end{interviewq}
